% Generated from committed aggregates; insert inside the relevant report section.
\subsection{Study design and inputs}

We built a research prototype for next-day eligible satellite fire detection in Blue Mountains National Park. The frozen model and its two development-selected recall policies did not establish reliable warning performance. The deliverable is an auditable ML comparison, not a deployed warning system or a validated ignition-probability map.

\subsubsection*{Prediction unit and evidence}

Each row is one retained ERA5-Land grid cell and one UTC feature day; the target is at least one eligible VIIRS detection on the following UTC day. Continuing fire counts as positive. Non-detection is not proof of no physical fire. Weather is aggregated over the grid cell, while terrain and labels concern its park portion; the group data owner must retain source and geometry documentation. ERA5-Land is retrospective reanalysis \cite{mlERA5Land}.

\subsubsection*{Temporal separation}

\begin{center}\small\setlength{\tabcolsep}{3pt}
\begin{tabular}{p{0.198\linewidth}p{0.525\linewidth}p{0.138\linewidth}}
\hline
Partition & UTC feature dates & Rows \\ \hline
Training & 1 July to 8 December 2019 & 6923 \\
Gap & 9 to 15 December 2019 & 301 \\
Validation & 16 to 24 December 2019 & 387 \\
Gap & 25 to 31 December 2019 & 301 \\
Final holdout & 1 January to 29 June 2020 & 7,783 \\
\hline\end{tabular}\end{center}

All cells from a date stayed together. The supplier's future-test boundary was preserved, with seven-day gaps rather than interleaved weeks \cite{mlTemporalEvaluation}. At least 10\% park overlap retained 43 cells. Rainfall warm-up excluded the first two feature days from fitting, leaving 6,837 training rows and 143 positive targets. Gaps reduce adjacency, not dependence between long-running fires.

\subsubsection*{Reviewed inputs}

The core has sixteen weather/terrain variables; the second version adds only the audited same-day fire\_today flag. Target-day counts, labels, dates, grid identifiers and coordinates are excluded as predictors. Long rainfall/dry-history and other current-fire candidates were not enabled because their warm-up or semantics needed further treatment. Fuel-moisture imagery and NDVI were not implemented.

\subsection{Models and evaluation procedure}

\subsubsection*{Matched models and stopping}

Training prevalence and persistence were explicit references. Logistic regression and XGBoost were each compared with sixteen and seventeen inputs on matched fitting/evaluation rows. Logistic scaling and balanced weights used training rows only. Tree fitting used training negatives divided by positives for class emphasis; calibration did not rebalance labels. XGBoost tests nonlinear thresholds/interactions \cite{mlChen2016}, rather than guaranteeing better performance.

\begin{center}\small\setlength{\tabcolsep}{3pt}
\begin{tabular}{p{0.490\linewidth}p{0.370\linewidth}}
\hline
XGBoost setting & Recorded value \\ \hline
Maximum depth & 3 \\
Learning rate & 0.05 \\
L2 leaf penalty & 5.0 \\
Row and feature subsampling & 0.8 and 0.9 \\
Early stopping & 25 rounds without AP improvement \\
Final 16 and 17 input tree counts & 6 and 8 \\
\hline\end{tabular}\end{center}

Average precision (AP) was computed from raw scores using sklearn, and explicitly maximised for stopping \cite{mlAveragePrecision}. The initial validation selected tree length and was development evidence. Later forward checks selected length on past inner data and refitted that fixed count before checking later dates. Missing-class windows remained visible; the dependent windows were not treated as independent replication.

\subsubsection*{Probability and reminder decisions}

A nondecreasing regularised sigmoid was fitted on base-model-excluded calibration rows, then checked on later dates, with seven-day gaps and no post-calibration base refit. The week-long design worsened later probability quality; the improving recent design had only one calibration and one assessment day. Calibration was therefore not adopted \cite{mlCalibration}. The highest calibration-only cutoff meeting each requested recall constraint was retained, with ties together.

\begin{center}\small\setlength{\tabcolsep}{3pt}
\begin{tabular}{p{0.267\linewidth}p{0.421\linewidth}p{0.172\linewidth}}
\hline
Development target & Exact raw cutoff & Comparison \\ \hline
80.0\% & 0.6658568382263184 & >= \\
90.0\% & 0.4193871319293976 & >= \\
\hline\end{tabular}\end{center}

The complete source, features, artifacts, thresholds and reporting rules were committed in freeze ef45cdf before one final release \cite{mlFinalEvidence}. No final fitting, calibration or threshold search occurred. AP was primary; confusion counts, probability quality and alert burden were reported overall, without current detections, and for every calendar month.
